\section{Outer index groups and systematic square codes}\label{sec:outer-codes}

The outer code has two roles.  Its systematic coordinates carry the data
on which the representation computes, and its redundant coordinates make
local consistency imply proximity in the metric that counts an entire
stored symbol as one position.  We construct the code, including its
matrix twists and simultaneous systematic encoders.  All vector-valued
statements below use joint support: a coordinate is erroneous when any
component is erroneous.  Their constants do not depend on the number of
components or stored slots.  The notation in this section is local to the
outer construction; in particular, the group constructed here is not the
affine group used later for the inner construction.

The index group must provide commuting left and right moves, constant
expansion, and only polylogarithmically many frequencies in each relevant
orbit. The first two properties support local consistency; the orbit
bound makes the local filling in Lemma~\ref{lem:outer-fill} possible.
The code and correction guarantees are summarized in
Theorem~\ref{thm:outer-codes} and Lemma~\ref{lem:outer-uncertainty}.

\subsection{A group of quasilinear order}

Let $h=2^\ell=O(N)$ be the padded tape length supplied by
Lemma~\ref{lem:macrograph}.  Fix a sufficiently large absolute integer
$C$, and choose an even $r$ with
\begin{equation}\label{oc:outer-size}
 \ell+C\le r\le\ell+C+2,
 \qquad M_o=2^r,\qquad G_0=\F_2^r.
\end{equation}
We increase $C$ whenever a fixed lower cutoff is needed.  Thus $4h\le M_o$,
$M_o=O(N)$, and $r=O(\log(N+2))$.  The bits of $x\in G_0$ are ordered
from least to most significant; $\underline{x}\in[0,M_o)$ denotes the
integer they represent.

Helfgott's ternary Goldbach theorem says that every odd integer greater
than five is a sum of three primes \cite[Main Theorem]{Helfgott2013}.
For our even $r$ we can therefore find primes
$p_1,p_2,p_3$ with $r-3=p_1+p_2+p_3$, by exhaustive search in $[r]$.
They need not be distinct or comparable in size.  Partition the $r$ bit
positions into blocks of lengths $p_a+1$, and let
\begin{equation}\label{oc:group}
 H=\prod_{a=1}^3\operatorname{PGL}_2(\F_{p_a}),
 \qquad \Gamma=G_0\rtimes H.
\end{equation}
Each factor acts on its projective line, hence permutes its block of bit
positions.  Since
$|\operatorname{PGL}_2(\F_q)|=q(q^2-1)$, we have
\begin{equation}\label{oc:group-order}
 |H|\le r^9,\qquad |\Gamma|=M_o|H|\le M_or^9.
\end{equation}
A generator of $\F_{p_a^2}^{\times}$, acting on its one-dimensional
$\F_{p_a}$-subspaces, induces a cycle $R_a$ of length $p_a+1$:
the scalar subgroup has order $p_a-1$ and its quotient acts freely and
transitively on these subspaces.  Label each block in this cycle order,
and denote its first unit vector by $e_a$.  The three finite fields,
permutation groups, cycles, and their tables have size polynomial in $r$
and can be found by exhaustive finite-field computations of that size.

The product is $(g,u)(w,v)=(g+uw,uv)$.  It is useful to write a group
element $k=(g,u)$ instead as $(x,u)$, where $x=u^{-1}g$.
Left and right multiplication by $(w,v)$ then have the formulas
\begin{equation}\label{oc:actions}
 (x,u)\longmapsto (x+u^{-1}v^{-1}w,vu),
 \qquad
 (x,u)\longmapsto (v^{-1}(x+w),uv),
\end{equation}
respectively.  These formulas also give polynomial-in-$r$ Boolean
algorithms for multiplication and inversion, using the small tables for
$H$ rather than a table for $\Gamma$.

\begin{lemma}[Uniform expansion]\label{oc:expansion}
For every fixed $0<\lambda<1$, the groups $\Gamma$ in
\eqref{oc:group} admit symmetric labelled Cayley multisets of a fixed
degree whose averaging operator has norm at most $\lambda$ on mean-zero
functions.  The same assertion holds for $H$.  Their labels can be
constructed deterministically in time polynomial in $N$ and represented
with $\operatorname{polylog}(N)$ bits of public data.  Multiplicities and identity
labels are permitted.
\end{lemma}

\begin{proof}
First consider a factor of $H$.  The reductions of the upper and lower
unit elementary matrices, with their inverses, have a uniform Cayley
gap in $\operatorname{SL}_2(\F_p)$.  This is the elementary-generator
case of the congruence gap discussed in \cite[p.~626]{BG2008}, and also
follows directly from \cite[Theorem~1]{BG2008}, since the corresponding
fixed subgroup of $\operatorname{SL}_2(\mathbb Z)$ is non-elementary.
Passing to the quotient by the center preserves the upper gap.  A
half-lazy walk has a uniform absolute gap.  The finitely many small
primes are absorbed by decreasing this fixed gap.

For odd $p$, the subgroup $\operatorname{PSL}_2(\F_p)$ is normal of
index two in $\operatorname{PGL}_2(\F_p)$.  Add a representative of the
nonsquare determinant class and its inverse; in the index-one case pad
with identities.  To check a uniform gap, work in an arbitrary unitary
representation.  The preceding subgroup gap bounds the distance of a
vector to its subgroup-invariant projection by a fixed multiple of its
maximum displacement under the subgroup labels.  Normality makes that
invariant subspace invariant under the extra label, which controls its
quotient of order at most two.  Thus distance to the full invariant
subspace is bounded by a fixed multiple of the maximum displacement
under all these labels.  The three factor projections commute; applying
this argument to the three factors gives the same statement for $H$,
with a fixed-size generating list.  The regular-representation gap used
here applies to every representation of the finite group, by its
decomposition into irreducible representations.

We next give a bounded list of translations which, together with these
$H$ labels, expands $\Gamma$.  Consider a block of $n'$ bits with cycle
$R$.  For a nonzero frequency $y\in\F_2^{n'}$, let $d'$ be the dimension
of the span of its cyclic shifts.  Identifying vectors with polynomials
modulo $Z^{n'}-1$, this span is a cyclic ideal.  Every such ideal of
dimension $d'$ is specified by a complementary monic divisor of degree
$d'$.  There are at most $2^{d'}$ such divisors and at most $2^{d'}$
vectors in each ideal.  In particular, there are at most $2^{2d'}$
frequencies with this span dimension.  This argument uses only the
ideal structure of the polynomial quotient and remains valid when
$Z^{n'}-1$ has repeated factors in characteristic two.

For a uniform seed $a\in\F_2^{n'}$, the correlation vector
$(y\cdot R^ta)_{t=0}^{n'-1}$ is uniform on a shift-invariant linear
space of dimension $d'$.  Take $j_0$ independent seeds.  In the product
of these spaces, the subset with total weight less than $j_0n'/4$ is
invariant under cyclic shifts and permutation of the seeds.  If this
subset is nonempty, its uniform distribution has the same marginal at
every coordinate, with probability of one at most $1/4$.  An information
set of $j_0d'$ coordinates determines a vector in the product space.
Entropy subadditivity consequently bounds the size of this subset by
$2^{j_0d'\mathcal H(1/4)}$, where $\mathcal H$ is binary entropy.
The subset of weight greater than $3j_0n'/4$ has the same bound, by the
marginal entropy estimate with probabilities at least $3/4$; it is
unnecessary for the space to be closed under complementation.
Thus a fixed sufficiently large $j_0$ makes the probability of any
failure at any nonzero frequency at most
\begin{equation}\label{oc:seed-bound}
 2\sum_{d'\ge1}
   2^{-\{j_0(1-\mathcal H(1/4))-2\}d'}<1.
\end{equation}
There exist $j_0$ seeds whose cyclic translates have absolute Walsh
bias at most $1/2$ at every nonzero frequency.  Enumerating all seed
tuples and checking their biases constructs them in time $2^{O(r)}$,
which is polynomial in $N$.

Take the Cartesian product of the three block seed lists.  Averaging
its translations under $\langle R_1\rangle\times\langle R_2\rangle
\times\langle R_3\rangle$ has absolute Walsh bias at most $1/2$ off
zero: the three factors multiply, and a nonzero frequency has a nonzero
block.  Averaging further over cosets in $H$ preserves this bound.
Denote by $T$ the average of the resulting $H$-conjugates of seed
translations.  On the orthogonal complement of $G_0$-invariants, $T$
has all eigenvalues in $[-1/2,1/2]$, by Fourier decomposition for $G_0$.

Let $v$ be a unit vector in a representation of $\Gamma$, and let
$\eta$ be its maximum displacement under the fixed $H$ labels and the
fixed product-seed translations.  The $H$-invariant projection $v_H$
satisfies $\|v-v_H\|\le K\eta$ for an absolute $K$.  Its displacement
under any seed translation is at most $(1+2K)\eta$.  Since $v_H$ is
$H$-invariant, the same holds for every $H$-conjugate of a seed.
Writing $P_{G_0}$ for projection to $G_0$-invariants gives
\[
 \|v_H-P_{G_0}v_H\|^2
 \le 2\langle v_H,(I-T)v_H\rangle
 \le (1+2K)^2\eta^2.
\]
Normality of $G_0$ implies that $P_{G_0}v_H$ is also $H$-invariant.
Distance to the $\Gamma$-invariant subspace is therefore at most a
fixed multiple of $\eta$.  The symmetric half-lazy average on our
fixed generating list has a uniform positive gap: its Dirichlet form
is the average of squared label displacements divided by a fixed
constant.  Its spectrum is nonnegative.  A fixed power of this walk
has mean-zero norm at most the requested $\lambda$.  Reversing words
shows that the powered list remains symmetric.  Its length is fixed,
and the same conclusion holds for left and right regular actions.
\end{proof}

\subsection{Common scalar local codes}

The distance and agreement estimates in this subsection hold over an
arbitrary field $J$, conditional on the scalar check matrices satisfying
the stated minor conditions.  We will choose the special outer field
and those matrices after fixing the constants in these estimates.

The local agreement constant must be fixed before the outer Cayley
degree is chosen.  To obtain this order, we make the label positions
the edges of a separate expander.  For every integer $s_b\ge1$, consider
the Margulis--Gabber--Galil multigraph~\cite{Margulis1973,GabberGalil1981}
on $\mathbb Z_{s_b}^2$ with the
eight labelled moves
\[
 (x\mathbin\pm2y,y),\quad (x\mathbin\pm(2y+1),y),\quad
 (x,y\mathbin\pm2x),\quad (x,y\mathbin\pm(2x+1)).
\]
Its normalized mean-zero norm is at most $5\sqrt2/8$ for every positive
integer modulus, including composite moduli.  The primary formal
development \cite[\texttt{mgg\_graph\_step},
\texttt{mgg\_numerical\_radius}]{MGGAFP} proves exactly this labelled,
symmetric, all-modulus statement.  Add eight identity labels for
half-laziness and take a sufficiently large fixed power.  Viewing the
result as a bipartite graph between two copies of $\mathbb Z_{s_b}^2$
gives a graph $B_0$ with fixed degree $d_0$, permutation labels
$1,\ldots,d_0$, and mean-zero cross-part norm
\begin{equation}\label{oc:base-parameters}
 \lambda_0<10^{-4},\qquad
 c_0=\lfloor d_0/16\rfloor,\qquad \delta_0=1/16.
\end{equation}
Take the fixed power large enough that $c_0\ge1$.
Let $n_b=s_b^2$ and $d=d_0n_b$.  Each of the $d$ edge coordinates
is named by its left base vertex and permutation label.  Both ends
of an outer edge use these same coordinate names.

For each of the four types $C_A^0,C_A^1,C_B^0,C_B^1$, choose two
$c_0\times d_0$ scalar matrices, one for each base side.  Impose the
corresponding checks at every base vertex, with incident edges in
permutation-label order.  There are eight check matrices in total;
their entries are shared by all stars of the same type.  We require
every $c_0$-column minor to be nonzero.  Section~\ref{oc:encoders-subsec}
constructs matrices meeting these conditions and the global encoder
conditions simultaneously.  Until then the following conclusions hold
for every choice satisfying the stated minors.  For a finite-dimensional
$J$-space $U$, notation such as $C_A^s(U)$ means that the scalar
constraints are applied componentwise to a word in $U^d$.

\begin{lemma}[Local distance and tensor agreement]\label{oc:local-agreement}
Set
\begin{equation}\label{oc:delta-one}
 \delta_1=\delta_0(\delta_0-\lambda_0)>0.
\end{equation}
Every nonzero word of any of the four local codes, over any
finite-dimensional $J$-space $U$, has at least $\delta_1d$ nonzero
coordinates.  There is a constant $C_0<\infty$, depending only on
$d_0,\delta_0,\lambda_0$, with the following property for every $s_b$,
every field and check matrices as above, every $U$, and every $s,t$.
If $F$ is a $d\times d$ table with rows in $C_B^t(U)$ and $D$ is such
a table with columns in $C_A^s(U)$, and their joint cell-disagreement
fraction is $z$, there is a table $W$ whose rows and columns obey both
codes and such that
\begin{equation}\label{oc:tensor-agreement}
 \Pr_i[W(i,\cdot)\ne F(i,\cdot)]
 +\Pr_j[W(\cdot,j)\ne D(\cdot,j)]\le C_0z.
\end{equation}
In particular, $C_0$ is independent of $d$, $|J|$, and $\dim_JU$.
\end{lemma}

\begin{proof}
The maximal-minor condition gives star distance at least
$c_0+1\ge\delta_0d_0$: a nonzero word of smaller support would give a
linear dependence among at most $c_0$ columns, which could be enlarged
to a dependence among $c_0$ columns.  This also proves the assertion
for vector entries, either by applying a nonzero scalar component or
by solving the independent-column equations componentwise.

For a nonzero whole local word, let $l,u$ be the fractions of left and
right base vertices incident to its support, and let $a=\max(l,u)$.
Its edge-support fraction $e$ obeys, by star distance and mixing,
\[
 \delta_0a\le e\le lu+\lambda_0\sqrt{lu}
 \le a(a+\lambda_0).
\]
Hence $a\ge\delta_0-\lambda_0$ and $e\ge\delta_1$.  Mixing here and
below follows by subtracting the two constant means and applying the
cross-part operator norm to the indicator functions.

For agreement, apply the A-base checks to the columns of $F$.  Mark
every base star at which any check syndrome is nonzero.  Each nonzero
syndrome, as a row indexed by the $d$ columns, is a scalar linear
combination of B-code rows, and is therefore a nonzero B-codeword over
$U$.  Its support has size at least $\delta_1d$.  In each such column
there must be a discrepancy with $D$ on an edge incident to the marked
star, since $D$ satisfies that star's check.  A discrepancy cell is
incident to at most two base stars.  If $w_0$ is the total marked-star
count divided by $n_b$ (counting both base sides), then
\begin{equation}\label{oc:marked-size}
 w_0\le \frac{2zd^2}{\delta_1dn_b}
       =\frac{2d_0z}{\delta_1}.
\end{equation}

We now confine every low-discrepancy difference column to edges inside
a small set of base vertices.  Deleting those row indices, together
with the columns having large discrepancy, will leave a rectangle on
which $F$ and $D$ agree.

Close the marked set by adding a vertex whenever at least
$(\delta_0/2)d_0$ incident edges lead to the set.  If the initial set
is empty it stays empty.  Otherwise suppose its normalized size first
reaches $w\ge2w_0$.  A single addition is at most $1/n_b\le w_0$, so
$w\le3w_0$.  Edges made internal by additions give internal edge
fraction at least $(\delta_0/2)(w-w_0)$.  Mixing gives the upper bound
$w^2/4+\lambda_0w/2$.  Since $w\ge2w_0$, these inequalities imply
$\delta_0\le w+2\lambda_0$.  They are incompatible when
$3w_0<\delta_0-2\lambda_0$.  For sufficiently small $z$, therefore,
the final closure $T$ satisfies $|T|/n_b\le2w_0$, with strict
inequality when $w_0>0$.
Every star outside $T$ still satisfies all checks on $F$, and fewer
than $(\delta_0/2)d_0$ of its edges enter $T$.

Fix a constant
\[
 0<z_0<\delta_0(\delta_0/2-\lambda_0).
\]
Consider a column whose discrepancy fraction is less than $z_0$.
The difference between its $F$ and $D$ columns satisfies all checks
outside $T$.  An outside vertex incident to its support consequently
has at least $\delta_0d_0$ support edges in total and more than
$(\delta_0/2)d_0$ support edges to other outside vertices.  The left
and right fractions of such outside vertices are each at most
$z_0/\delta_0$, by counting their total support edges.  If $a$ is
their maximum fraction, mixing on the outside support would give
\[
 (\delta_0/2)a\le a(a+\lambda_0),
 \qquad a\le z_0/\delta_0<\delta_0/2-\lambda_0.
\]
Thus $a=0$.  The difference column vanishes on every edge having an
endpoint outside $T$.

Delete the row indices corresponding to edges internal to $T$, and
delete columns with discrepancy fraction at least $z_0$.  The former
fraction is at most $|T|/(2n_b)\le w_0$, since each internal edge consumes
two incidences in $T$, and the latter fraction is at most $z/z_0$.
The two tables agree on the complementary rectangle.  Choose a fixed
$z_*>0$, depending only on the base parameters, so small that for
$z<z_*$ the closure inequality holds and each deleted fraction is
less than $\delta_1$.  Puncturing either scalar local code to its
retained coordinates is then injective by distance.  Its inverse
on the image is a scalar linear map.  Lift the common rectangle in
the B direction and then in the A direction.  Applying scalar linear
maps in one direction preserves the constraints in the other; hence
the resulting table belongs to the full tensor code.  By injectivity
it agrees with every retained whole row of $F$ and every retained
whole column of $D$.  This gives \eqref{oc:tensor-agreement} with
$2d_0/\delta_1+1/z_0$ for $z<z_*$.  For $z\ge z_*$, use the zero
tensor table and the bound $2\le(2/z_*)z$.  The maximum of these two
constants is a suitable $C_0$.  All lifting maps and all support
counts act on joint vector entries, so no dimension factor occurs.
\end{proof}

We now fix, once and for all,
\begin{equation}\label{oc:lambda-choice}
 0<\lambda\le\min\{(8C_0)^{-1},\delta_1/16\}.
\end{equation}
Apply Lemma~\ref{oc:expansion} to obtain an outer Cayley list of degree
$D$.  Choose $s_b$ to be a positive multiple of $D$.  Then
$d=d_0s_b^2$ is divisible by $D$, so label the $d$ local coordinates
by a balanced repetition of the outer list.  This repetition changes
neither averaging operator nor its norm.  Use such lists
$(a_i)_{i=1}^d$ and $(b_j)_{j=1}^d$ for the A and B directions.
All of $d_0,c_0,\lambda_0,\delta_1,C_0,\lambda,s_b,d$ are now fixed
constants.  We now specialize $J$; the field choice will not change
these constants.

Choose an odd characteristic independently of the primes $p_a$
and independently of the characteristic-two translation group.  Put
$m=M_o$.  Linnik's theorem, for example the effective bound in
\cite[Theorem~1.1]{Xylouris2011}, supplies a prime
\begin{equation}\label{oc:prime}
 p\equiv1\pmod m,\qquad m<p\le C_Lm^{L}
\end{equation}
for absolute constants $C_L,L$.  Searching integers and testing
primality, even by trial division, takes polynomial time in $m$.
Search the prime field for an element $\zeta$ of exact order $m$;
such an element exists because its multiplicative group is cyclic.
Since $m$ is a power of two, testing $\zeta^m=1$ and
$\zeta^{m/2}\ne1$ suffices.  We will use
\begin{equation}\label{oc:field}
 J=\F_{p^4}.
\end{equation}
The proof below verifies that degree four supplies enough room for
every simultaneous outer-code choice.  One can find an irreducible
quartic by enumeration and represent $J$ as its polynomial quotient.
Field operations and the root $\zeta$ then use $O(\log N)$ bits and
polynomial-in-$\log N$ Boolean computations.  Our fixed lower cutoff
also ensures $p>m>|H|$, so both $m$ and $|H|$ are invertible in $J$.

\subsection{Twists and simultaneous systematic extension}
\label{oc:encoders-subsec}

There are four outer parts
\[
 V=\Gamma\times\{0,1\}\times\{0,1\}.
\]
At $(k,s,t)$ a full slot stores a $d\times d$ matrix $W_k^{st}$ over
a finite-dimensional $J$-space $U$, with columns in $C_A^s(U)$ and
rows in $C_B^t(U)$.  The A and B edges impose, respectively,
\begin{equation}\label{oc:edge-equations}
 W_{a_ik}^{1t}(i,j)=\alpha_i W_k^{0t}(i,j),
 \qquad
 W_{kb_j}^{s1}(i,j)=\beta_j W_k^{s0}(i,j).
\end{equation}
The maps $\alpha_i,\beta_j$ are invertible $J$-linear maps on $U$,
independent of $k,s,t$, and satisfy
$\alpha_i\beta_j=\beta_j\alpha_i$ for every $i,j$.  An A edge of
label $i$ transports the entire B-coded row $i$, while a B edge of
label $j$ transports the entire A-coded column $j$.  A reverse edge
uses the inverse map.  At their intersection the two transports give
\[
\begin{array}{ccc}
 W_k^{00}(i,j)&\xrightarrow{\ \alpha_i\ }&W_{a_i k}^{10}(i,j)\\
 {\scriptstyle\beta_j}\big\downarrow&&
       \big\downarrow{\scriptstyle\beta_j}\\
 W_{k b_j}^{01}(i,j)&\xrightarrow{\ \alpha_i\ }&W_{a_i k b_j}^{11}(i,j).
\end{array}
\]
Left and right multiplication on $\Gamma$ commute, so the two routes
reach the same index; commuting entry maps give the same entry value.

The scalar type has $U=J$ and identity twists.  The programs in
Section~\ref{sec:outer-programs} will also need to multiply scalar
data $v:\Gamma\to J$ by $(-1)^{e\cdot x}$ for a distinguished bit
$e=e_a$.  The matrix types below supply edge equations compatible
with this operation.  The programs will use them to retain one axis
of code constraints while changing the other.
For each of the three distinguished bits $e=e_a$, define a
representation on the basis
$({\bf e}_{t'})_{t'\in H}$ by
\begin{equation}\label{oc:rho}
 \rho_e(g,u){\bf e}_{t'}
   =(-1)^{(ut'e)\cdot g}{\bf e}_{ut'}.
\end{equation}
Its entries are signed permutation entries in the prime field.
Indeed, for $k=(g,u)$ and $k'=(g',u')$, the exponent in
$\rho_e(k)\rho_e(k'){\bf e}_{t'}$ is
\[
 (u't'e)\cdot g'+(uu't'e)\cdot g
  =(uu't'e)\cdot(g+ug'),
\]
which proves multiplicativity.  In a matrix type let
$U=\operatorname{Mat}_{|H|}(J)$ and independently choose each axis
to be untwisted or to use
\begin{equation}\label{oc:matrix-twists}
 \alpha_i D=\rho_e(a_i)D,
 \qquad \beta_j D=D\rho_e(b_j).
\end{equation}
Left and right matrix multiplication commute.  Right multiplication
in \eqref{oc:matrix-twists} specifies edge maps; no convention about
its order as a group representation is needed.  This gives at most
$1+3\cdot4=13$ full types.  Finite tuples, with arbitrary repetitions
of these types, use direct sums of their entry spaces and block
diagonal twists.

On either local-label graph, let the parity set consist of all edges
with permutation number in $[4c_0]$.  Its size and complement are
\begin{equation}\label{oc:parity-set}
 q_{\rm par}=4c_0s_b^2,\qquad
 I_A=[d]\setminus\text{A-parity set},\qquad
 I_B=[d]\setminus\text{B-parity set}.
\end{equation}
Both systematic sets are nonempty; in fact they contain at least
$3d/4$ labels.

\begin{lemma}[Simultaneous systematic encoders]\label{oc:systematic}
One can choose the eight scalar check matrices over $J$ so that all
their maximal minors are nonzero and every A and B one-axis system
of the thirteen types above is systematically invertible.  Precisely,
arbitrary data at the systematic labels of part zero have a unique
extension to both parts satisfying that axis's local checks and edge
equations.  Consequently arbitrary data on
$\Gamma\times I_A\times I_B$ in part $(0,0)$ have a unique full
extension satisfying both axes, and the two orders of axis extension
give the same answer.  The matrices and the full encoders are
deterministically constructible in polynomial time in $N$.  These
assertions continue to hold for arbitrary finite direct sums of the
listed types, without any further field or check-matrix choices.
\end{lemma}

\begin{proof}
Consider one A-axis system.  Eliminate part-one variables with the
A-edge equations.  At each $k$, the two types of base-star checks in
part zero contribute $2c_0s_b^2$ equations; those in part one contribute
the same number after substitution.  There are exactly
$q_{\rm par}=4c_0s_b^2$ parity-label variables at part zero for each
$k$, so the resulting parity system is square, also after scalarizing
$U$.  Its entries are linear polynomials in the entries of the four
A check matrices.

Its determinant polynomial is nonzero in the prime characteristic.
For this purpose partition the $4c_0$ parity permutation numbers into
four disjoint sets of $c_0$ numbers, one for each pair consisting of an
outer part and a base side.  Specialize each associated check matrix
to selectors of its set.  For a fixed permutation number, every
label-index edge occurs exactly once at the specified base side.
At outer part zero its selected coordinate is unshifted; at part one
it is pulled back through an invertible group shift and an invertible
entry twist.  Thus every scalarized parity coordinate is selected
exactly once, up to an invertible block and a permutation.  The
specialized parity matrix is invertible.  The same proof works for
the B axis and for every allowed twist type.  It does not require
this selector specialization itself to have the desired maximal
minors: it separately witnesses that each determinant is a nonzero
polynomial.  Each maximal-minor polynomial is also nonzero.  Their
product is nonzero because a polynomial ring over a field is an
integral domain.

Here is a degree bound that justifies this field choice without
reference to any future number of slots.  There are exactly
\begin{equation}\label{oc:number-variables}
 8c_0d_0
\end{equation}
scalar indeterminates.  Scalarizing a one-axis parity matrix gives
dimension at most
\begin{equation}\label{oc:parity-dimension}
 q_{\rm par}|\Gamma|\dim_JU
 \le q_{\rm par}m|H|^3\le q_{\rm par}mr^{27}.
\end{equation}
Its determinant has degree at most this number.  At most
$2(1+3\cdot4)=26$ axis systems cover every required type, with
redundancies harmless.  The product of the maximal minors of the
eight check matrices has degree $8c_0\binom{d_0}{c_0}$.  Thus the
total degree is at most
\begin{equation}\label{oc:determinant-degree}
 26q_{\rm par}mr^{27}+8c_0\binom{d_0}{c_0}<m^3.
\end{equation}
For the last inequality increase the fixed cutoff so that
$r^{27}\le m$, $m\ge2\cdot26q_{\rm par}$, and
$m^2\ge2\cdot8c_0\binom{d_0}{c_0}$.  The first term is then at
most $m^3/2$ and the second at most $m^2/2<m^3/2$.
The exponent three is independent of $s_b$; only the cutoff depends
on the already fixed local degree.

A nonzero polynomial of total degree less than a finite field's
cardinality cannot vanish on every tuple over that field.  For
completeness, view it as a polynomial in the last variable; by
induction choose the other variables so that a nonzero leading
coefficient remains nonzero, and then use the univariate root bound.
Since $|J|=p^4>m^3$, such a tuple exists in $J$.  Enumerate all tuples
of the fixed number \eqref{oc:number-variables} of field elements,
testing maximal minors and parity ranks, until one succeeds.  The
field cardinality and every matrix dimension are polynomial in $N$,
so this is a deterministic polynomial-time algorithm with a fixed
exponent.  A direct sum of already invertible parity systems is
invertible with the direct sum of their inverses.  It introduces no
new polynomial condition, even when the number of slots grows with
$N$.

It remains to justify two-axis extension.  Regard an A parity matrix
as a matrix of operators on $U^\Gamma$: each entry is a scalar
combination of left translations and A twists.  B entries are scalar
combinations of right translations and B twists.  They commute
entrywise across axes.  An invertible A parity matrix $L_A$ therefore
intertwines the diagonal actions of any B-entry operator $T$ on its
domain and codomain.  From $L_AT=TL_A$ one obtains
$L_A^{-1}T=TL_A^{-1}$, with the same indexed domain and codomain
interpretation.  Thus the operators making up either axis encoder
commute with every entry operator of the other axis encoder, including
its inverse parity solve.

Start with the systematic rectangle at part $(0,0)$.  Extend along A
for every systematic B column and then along B for every resulting
A row.  The second extension preserves all first-axis constraints
by this commutation.  Both encoders leave their systematic entries
unchanged.  Reversing their order gives the same map.  Finally, one-axis
uniqueness applied in both directions shows that a full word is
uniquely determined by its systematic rectangle.  Gaussian elimination
computes these extensions in polynomial time in $N$, or in polynomial
time in the total explicit input size for an arbitrary direct sum.
\end{proof}

For $v:\Gamma\to U$, give every systematic cell in part $(0,0)$ the
same value $v(k)$.  Its full extension is denoted
\begin{equation}\label{oc:encoder}
 E_{\alpha,\beta}v,
\end{equation}
or $Ev$ when its type is specified.  We call this the
\emph{duplicated-payload} encoder.  It is linear and injective, and
its rate in field coordinates is exactly $1/(4d^2)$: its domain has
$|\Gamma|\dim_JU$ coordinates and its stored word has
$4d^2|\Gamma|\dim_JU$ coordinates.  In the largest individual type
the entry dimension is $|H|^2\le r^{18}$.  Therefore one full slot
occupies $\operatorname{polylog}(N)$ bits at each of $|\Gamma|=O(N\operatorname{polylog} N)$
sites.  Any fixed power of $\log(N+2)$ slots preserves this size form.

\begin{remark}[Operator invariance]\label{oc:orbit-invariance}
The preceding inverse argument also justifies restriction to frequency
orbits, as required by the local filling algorithms.  Expand a function
of $(x,u)$ in the Walsh characters $(-1)^{\nu\cdot x}$, using $m^{-1}$
for the inverse transform over $J$.  By \eqref{oc:actions}, a left pull
preserves each $x$-frequency, with a phase depending on $u$, whereas
a right pull permutes frequencies by $H$.  Constant entry twists do
not change the frequency.  For any $H$-orbit $O\subseteq G_0$, both
the $O$-frequency subspace and its complementary frequency subspace
are invariant under each parity operator.  Its domain and codomain
have equal dimensions on this subspace: the parity-label and
check-label counts are both $q_{\rm par}$.  Global invertibility
therefore implies invertibility of the restriction.  Its scalar
dimension is at most $q_{\rm par}|O||H|\dim_JU\le
q_{\rm par}|H|^4$, polynomial in $r$, so the restriction and its
inverse can be computed by polynomial-in-$\log N$ field operations.
This is an algebraic consequence of the explicit parity matrices;
it does not require reading unspecified global values of a word.
\end{remark}

\subsection{Proximity and one-axis gaps}

Write $\dist_V$ for relative Hamming distance on the four-part vertex
set $V$.  When all four parts at the same $k$ are packed into one
symbol, write $\dist_\Gamma$.  For any two words,
\begin{equation}\label{oc:packing-metric}
 \dist_V\le\dist_\Gamma\le4\dist_V.
\end{equation}
An \emph{A-only word} retains the A local constraints and A agreements
at all four parts; any other label indices and any collection of slots
are regarded as coordinates of a single row symbol.  No B validity
or B agreement is asserted.  Define B-only words symmetrically.

For a tuple of full slots a local check rejects when any component
fails its scalar local checks, and an edge is conflicting when any
slot fails its agreement on the whole row or column carried by that
edge.  In particular, a growing tuple does not enlarge the set of
address maps or cause separate per-slot error charges.

\begin{lemma}[Full-code proximity]\label{oc:full-proximity}
There are fixed $\theta_*>0$ and $K_*<\infty$ such that the following
holds for every finite tuple of allowed full types and every received
word.  Check local validity and all A and B agreements at a uniformly
random packed center, rejecting if any check there fails.  If the
rejection fraction is $\theta\le\theta_*$, the word is at packed
distance at most $K_*\theta$ from a word satisfying every local and
edge constraint.  The constants are independent of the tuple length,
the field size, and the entry dimensions.  Duplicated-payload
constraints are addressed separately below.
\end{lemma}

The proof uses the decreasing-disagreement method of
\cite[Section~4.2, Algorithm~1 and Propositions~4.7--4.8]{DEL+2022}.
We prove the version for joint tuples with commuting twists needed here.

\begin{proof}
First replace every locally invalid whole tuple at a vertex by zero.
There are at most $\theta|V|$ such vertices.  Initially at most
$d\theta|V|$ undirected edges conflict, and this replacement creates
at most $2d\theta|V|$ further conflicts.  Starting with this locally
valid word, repeatedly replace a whole vertex by a locally valid tuple
whenever that strictly decreases its number of incident conflicting
edges.  The total integer conflict count strictly decreases, so at
most $3d\theta|V|$ replacements occur.  This argument establishes
existence and termination; it is not asserted to be an efficient
search through the alphabet.  At most $(1+3d)\theta|V|$ vertices have
been touched, and the final conflict count is no larger than the
count before the replacements.

At a stable vertex $v$, let $p_A(v),p_B(v)$ be the fractions of its
A and B edges that conflict, and put
$p(v)=p_A(v)+p_B(v)$ and $\mu=\E_Vp$.  Neighbors across A edges
suggest rows, transported by the appropriate twist or inverse twist.
They are valid B-code rows.  Neighbors across B edges similarly
suggest valid A-code columns.  At cell $(i,j)$ these suggestions
agree unless at least one of the two opposite square edges conflicts:
the common far corner and commuting twists supply the equality when
both opposite edges agree.  If $P_A,P_B$ average over the respective
neighbor maps, their cell-disagreement fraction is at most
$P_Ap_B(v)+P_Bp_A(v)$.  Lemma~\ref{oc:local-agreement} gives a local
replacement whose incident conflict fraction is at most $C_0$ times
this number.  Stability therefore gives pointwise
\begin{equation}\label{oc:stability}
 p\le C_0(P_Ap_B+P_Bp_A)\le C_0(P_Ap+P_Bp).
\end{equation}

All norms and inner products in this proof use uniform normalized
counting measure on $V$.  Let $\bar p$ be the function equal to the
mean of $p$ on each of the four parts.  Nonnegativity gives
$\|\bar p\|_2\le2\mu$.  Each neighbor operator permutes the part
means and has norm at most $\lambda$ on their orthogonal complement.
Consequently \eqref{oc:stability} implies
\[
 \|p\|_2\le2C_0(2\mu+\lambda\|p\|_2),
 \qquad \|p\|_2\le C_1\mu,\quad C_1=8C_0,
\]
by \eqref{oc:lambda-choice}.  This bound has no dependence on $d$.

Fix an A label $i$, let $S_i$ be the vertices whose incident A edge
of that label conflicts, and let $p_i=\Pr_V(S_i)$.  At such an edge
the difference of its two transported rows is a nonzero B-codeword,
so it has at least $\delta_1d$ nonzero cells.  In each corresponding
B direction, either a B edge conflicts at one A endpoint or an
$i$-label A edge conflicts at the opposite B neighbor.  The
involution exchanging the two A endpoints preserves $S_i$.  Averaging
this implication, and then subtracting the four part means in the
mixing term, yields
\begin{align}
 \delta_1p_i
 &\le 2\E_V[1_{S_i}p_B]
      +\langle1_{S_i},P_B1_{S_i}\rangle \notag\\
 &\le 2\E_V[1_{S_i}p_B]+4p_i^2+\lambda p_i.
 \label{oc:conflict-propagation}
\end{align}
Indeed, the squared norm of the four part means of $1_{S_i}$ is at
most $4p_i^2$, and the squared norm of its remaining part is at most
$p_i$.  Also $p_i\le d\mu$.  For sufficiently small $\mu$, depending
on the now fixed $d$, we have $4d\mu+\lambda<\delta_1/2$.
Average \eqref{oc:conflict-propagation} over $i$ and repeat the
argument with A and B exchanged.  It follows that
\begin{equation}\label{oc:stable-zero}
 (\delta_1/2)\mu
 \le4\E_V[p_Ap_B]\le4\E_V[p^2]\le4C_1^2\mu^2.
\end{equation}
If $0<\mu<\delta_1/(8C_1^2)$ this is impossible.  Since the initial
zeroing left at most $3d\theta|V|$ conflicts, stability has
$\mu\le6\theta$.  For example, choose $\theta_*$ so that
$6\theta_*<\min\{\delta_1/(32d),\delta_1/(16C_1^2)\}$.
Then $\mu=0$.  The final word is exact, and
\eqref{oc:packing-metric} gives $K_*=4(1+3d)$ as a permissible
proximity constant for this normalization.  Common scalar local
checks, block diagonal twists, and joint conflict events justify
every step for an arbitrary tuple without a slot-count factor.
\end{proof}

\begin{lemma}[One-axis distance and exact comparisons]\label{oc:one-axis}
Let $\Delta_A=(\delta_1-\lambda)/4>0$.  Every nonzero A-only or
B-only word with the specified local constraints and invertible twists
has relative vertex support at least $\Delta_A$ on $V$.  The same
holds for differences of two words with the same axis constraints,
and for arbitrary finite joint tuples.  It also holds for full words.

Among exact full words, the duplicated-payload condition has a fixed
gap under the following checks: at $s=0$, compare all systematic A
rows as whole rows, at both values of $t$; at $t=0$, compare all
systematic B columns as whole columns, at both values of $s$.
If any payload is not duplicated, the packed-center rejection
fraction is at least $\Delta_A$.  The tests accept every word $Ev$.

More generally, suppose an expression $X$ and a full target $Y=Ev$
both obey the same A constraints, and the systematic rectangle of
$X$ is duplicated with payload $w$.  Equality of their whole
systematic B columns at $t=0$, including both $s$ parts, holds if
and only if $v=w$.  If $v\ne w$, the comparison rejects at least a
$\Delta_A$ fraction of packed centers.  The symmetric assertion
holds for a shared B axis and comparisons of whole systematic A
rows at $s=0$.  No constraint on the other axis of $X$ is needed.
\end{lemma}

\begin{proof}
For one A slice with fixed $t$, a nonzero vertex has at least
$\delta_1d$ nonzero row symbols by local distance.  Each is carried
to a nonzero symbol at its opposite endpoint, because twists are
invertible.  If $l,u$ are the two vertex-support fractions and
$a=\max(l,u)$, the fraction $e$ of nonzero carried edges obeys
\[
 \delta_1a\le e\le lu+\lambda\sqrt{lu}\le a(a+\lambda).
\]
Thus $a\ge\delta_1-\lambda$.  Support on these two parts is at
least $(\delta_1-\lambda)/2$ in their own normalization and at least
$\Delta_A$ in the four-part normalization.  At least one slice is
nonzero.  The argument for B, differences, and joint entries is
identical.

For two systematic A row labels, their difference at $s=0$, considered
over both $t$ parts and all $k$, is a B-only word: the scalar B checks
and B twists are the same for both rows.  If nonzero it has the
distance just proved.  Differences between systematic B columns are
A-only words.  A nonduplicated systematic rectangle contains two
unequal entries connected either by a change of its A label or a
change of its B label, so one of these comparisons is nonzero.
Conversely, for a duplicated payload, B-axis uniqueness extends the
same systematic row data identically at every systematic A label,
and A-axis uniqueness does the analogous thing for columns.

For the last assertion, if $v=w$, A-axis uniqueness identifies the
two whole columns at every systematic B label.  If $v\ne w$, a
systematic payload cell differs somewhere, so one of the compared
column differences is a nonzero A-only word and has support at least
$\Delta_A$.  The B-sharing assertion is proved in the same way.
These comparisons use whole one-axis codewords, which is why an
arbitrary localized payload disagreement has a constant exact gap.
\end{proof}

The next estimate will control a growing tuple of A-only scratch
slots against one exact reference.  It uses the tuple's joint edge
conflicts, with no factor for its number of components.

\begin{lemma}[One-axis near-code estimate]\label{oc:near-code}
Let $W$ and $V_0$ be locally A-valid joint tuples on $V$, and suppose
$V_0$ has exact A agreement with the same twists as $W$.  Let
$p=\dist_V(W,V_0)$ and let $z_A$ be the A-edge conflict fraction of
$W$, normalized as the probability that a uniform vertex and a
uniform A label give a conflict.  Then
\begin{equation}\label{oc:near-code-inequality}
 \delta_1p\le z_A+4p^2+\lambda p.
\end{equation}
In particular, if $p\le\delta_1/16$, then
$p\le2z_A/\delta_1$.  Both statements hold for every finite tuple
of A-only and full slots, uniformly in its size.  They also hold
with A and B exchanged.
\end{lemma}

\begin{proof}
Let $S$ be the joint vertex-disagreement set.  At each vertex in $S$,
local distance gives at least $\delta_1d$ nonzero difference rows.
For each corresponding label, either the other endpoint also lies
in $S$, or agreement of $V_0$ forces an actual conflict of $W$.
Consequently
\[
 \delta_1p\le z_A+\langle1_S,P_A1_S\rangle
 \le z_A+4p^2+\lambda p,
\]
by the four-part mixing bound used above.  For
$p\le\delta_1/16$ and \eqref{oc:lambda-choice},
$4p+\lambda<\delta_1/2$, giving the stated rearrangement.
The proof concerns one fixed received word and one joint reference.
It can be applied to successively larger prefixes of a tuple without
modifying the received word or accumulating a number-of-prefixes
loss.
\end{proof}

We also allow a clock or control tuple repeated identically at every
$k$.  For any received tuple, let its distinct-value class proportions
be $u_j$.  Equality across an expanding symmetric A list has probability
at most
\begin{equation}\label{oc:clock-gap}
 \sum_j u_j^2+\lambda\left(1-\sum_j u_j^2\right).
\end{equation}
This follows by applying the mean-zero bound to each class indicator
and summing.  If the equality rejection is $\eta$, then
$1-\max_j u_j\le1-\sum_j u_j^2\le\eta/(1-\lambda)$.
Thus a small rejection fraction implies proximity to one constant
tuple, irrespective of its bit length.

Use unique bit representations of field elements and canonical zero
padding.  At a packed center check parsing, ranges, and padding.
Replace a malformed symbol by a fixed parsable default when reasoning
from acceptance.  If all checks read from one fixed finite list
$\Pi$ of permutation addresses, modifying a set of $u|\Gamma|$
symbols can change their center outcomes on at most
$|\Pi|u|\Gamma|$ centers, by taking the union of inverse images.
This applies even when many slots are checked together or branches
use only subsets of the list.  Combining this observation with
Lemma~\ref{oc:full-proximity}, \eqref{oc:clock-gap}, and the
duplicated-payload gap of Lemma~\ref{oc:one-axis} gives proximity
to exact duplicated full encodings and a repeated clock whenever
the total packed rejection is sufficiently small.  Indeed, the
one-time repair changes only $O(\theta)$ centers and hence leaves
duplication rejection $O(\theta)$; choose its threshold below the
fixed exact gap.  Any required public range for a corrected clock
is likewise enforced, since a constant illegal value fails its
local range test everywhere.  These constants may depend on the
fixed address-list size, but not on the slot count.

\begin{theorem}[Outer-code interface]\label{thm:outer-codes}
There exist fixed constants $d$, $\delta_1>0$, $C_0<\infty$,
$\lambda>0$, $\theta_{\mathrm{out}}>0$, and $K_{\mathrm{out}}<\infty$, with
$\lambda\le\min\{(8C_0)^{-1},\delta_1/16\}$, having the following
uniform properties for every input length $N$ and its padded tape
length $h=O(N)$, after the fixed padding in \eqref{oc:outer-size}.
\begin{enumerate}[label=(\roman*)]
\item The group $\Gamma=\F_2^r\rtimes H$ has
$|\Gamma|\le M_or^9=O(N\operatorname{polylog} N)$, where $M_o=2^r=O(N)$.
Its A and B label lists are symmetric, of degree $d$, and have
mean-zero norm at most $\lambda$.  There is a field
$J=\F_{p^4}$ with $p>M_o>|H|$, $p\equiv1\pmod{M_o}$, and a
primitive $M_o$-th root $\zeta\in\F_p$.  All public group and field
data and the common scalar check matrices are deterministically
constructible in polynomial time in $N$, and the data needed at an
index have $\operatorname{polylog}(N)$ size and Boolean computation cost.
\item Each of the scalar and three matrix-representation families,
with either axis independently twisted or untwisted as specified
in \eqref{oc:matrix-twists}, has the simultaneous systematic
encoder \eqref{oc:encoder}.  Arbitrary systematic rectangle data
extend uniquely; the A and B extensions commute.  Duplicated payload
encoding is injective, linear, and has rate $1/(4d^2)$ in field
coordinates.  A full slot has $\operatorname{polylog}(N)$ bits per packed site.
Arbitrary finite tuples reuse the same checks and inverses.
\item The local distance and tensor-agreement statement of
Lemma~\ref{oc:local-agreement} holds uniformly in field size and
joint vector dimension.  Local validity, edge agreement, and
duplicated-payload comparisons have perfect completeness.  For
every finite tuple of full slots, if their joint packed-center
rejection fraction is $\theta\le\theta_{\mathrm{out}}$, their packed
distance to a tuple of exact duplicated encodings is at most
$K_{\mathrm{out}}\theta$.  Canonical parsing and a repeated control
tuple can be included with the same form of conclusion, changing
only the fixed constants.
\item The one-axis distance and shared-axis comparison gaps of
Lemma~\ref{oc:one-axis} hold with $\Delta_A=(\delta_1-\lambda)/4$.
For every actual and reference joint tuple satisfying the hypotheses
of Lemma~\ref{oc:near-code},
$\delta_1p\le z_A+4p^2+\lambda p$ and, when
$p\le\delta_1/16$, $p\le2z_A/\delta_1$.  A-only slots require
no B constraints for these conclusions.  The symmetric B assertions
hold as well.
\end{enumerate}
All constants are independent of $N$ and of the number of later
program slots.  The field search has only $8c_0d_0$ indeterminates,
at most $26$ parity determinant factors, and total degree less than
$M_o^3$; in particular it causes no feedback into $C_0$, $\lambda$,
or $d$.
\end{theorem}

\begin{proof}
Parts (i) and (ii) follow from Lemmas~\ref{oc:expansion} and
\ref{oc:systematic}, the field construction, and the coordinate
count after \eqref{oc:encoder}.  Part (iii) is
Lemma~\ref{oc:local-agreement}, Lemma~\ref{oc:full-proximity}, and
the parsing and exact-gap argument following \eqref{oc:clock-gap}.
Decrease $\theta_{\mathrm{out}}$ and increase $K_{\mathrm{out}}$ once to
cover those fixed costs.  Part (iv) is
Lemmas~\ref{oc:one-axis} and \ref{oc:near-code}.  The degree and
dependency assertions are \eqref{oc:number-variables}--
\eqref{oc:determinant-degree} and the parameter order fixed in
\eqref{oc:lambda-choice}.  No external local-testability theorem is
used for these specialized conclusions.
\end{proof}

\subsection{Correction with value-independent uncertainty supports}

The next guarantee is stronger than correction of a single fixed
word.  It tracks a set on which values are wholly unspecified, and
therefore applies to separate bounded-cone simulations with different
choices on that set.

\begin{lemma}[Outer uncertainty shrinkage]\label{lem:outer-uncertainty}
Fix the data in Theorem~\ref{thm:outer-codes}, and any finite tuple of
full or A-only slots with their allowed A twists.  Optionally include
a control tuple constant over $\Gamma$ and canonical padding.
There is a deterministic synchronous local correction map $\mathcal C$
with the following property.  Let $V_0$ be any reference word satisfying
all A local constraints and A agreements, with constant control and
canonical parsing.  Let $S\subseteq\Gamma$ and let $W$ be any input
equal to $V_0$ outside $S$, with arbitrary values or malformed symbols
on $S$.  Then $\mathcal C W=V_0$ outside
\begin{equation}\label{oc:uncertainty-set}
 S'=
 \left\{k:\frac1d\#\{i:a_ik\in S\}\ge\delta_1/2\right\}
 \ \cup\
 \left\{k:\frac1d\#\{i:a_i^{-1}k\in S\}\ge\delta_1/2\right\}.
\end{equation}
In particular, $S'$ depends on $S$ alone.  If $u=|S|/|\Gamma|<\delta_1/4$,
then
\begin{equation}\label{oc:uncertainty-shrink}
 \frac{|S'|}{|\Gamma|}
 \le\frac{2\lambda^2u}{(\delta_1/2-u)^2}\le u/2.
\end{equation}
The local map uses a fixed number of permutation addresses and time
polynomial in the symbol bit length and the public local-parameter
length, with total default behavior on arbitrary inputs.

For every fixed integer $q\ge0$, the $q$-round map has a bounded
input cone consisting of fixed permutations of its target; the bound
may depend on $q$ but not on $N$.  For an arbitrary initial $S$, the
recursively defined supports contain the errors simultaneously for
all completions on $S$.  If $|S|/|\Gamma|<\delta_1/4$, their densities
are at most $2^{-q}|S|/|\Gamma|$ after $q$ rounds.  Different target-cone
simulations may use different such completions, provided each simulation
assigns one consistent initial value to every address repeated within
that cone.  Under this small-initial-support hypothesis,
$O(\log|\Gamma|)$ rounds correct the whole word exactly in time
polynomial in its explicit size and $N$.
\end{lemma}

\begin{proof}
At a target part $(k,0,t)$, take the row of label $i$ from
$(a_ik,1,t)$ and apply $\alpha_i^{-1}$.  At $(k,1,t)$, take it from
$(a_i^{-1}k,0,t)$ and apply $\alpha_i$.  Combine all other indices
and all slots into one row-symbol vector.  Decode these $d$ row
suggestions in $C_A^s$ up to strictly fewer than $\delta_1d/2$
joint row errors.  There is at most one such local word by the
distance in Lemma~\ref{oc:local-agreement}.  A total deterministic
implementation enumerates all subsets $T\subseteq[d]$ with
$|T|>d-\delta_1d/2$, solves the local scalar constraints together
with the suggested values on $T$ componentwise, and verifies that
the candidate has fewer than $\delta_1d/2$ joint row errors.
Such a puncturing is injective by local distance, so each consistent
system gives at most one candidate.  If the distance promise holds,
the set of all correct suggestions is one enumerated $T$ and gives
the required word.  Return zero if no candidate passes.  Malformed
neighbor symbols are first replaced by the fixed parsable default.
Restore each control bit by majority on the forward A-neighbor list,
with fixed tie behavior, and write canonical padding.

If $k\notin S'$, each relevant list contains fewer than
$\delta_1d/2$ positions of $S$.  Every other row suggestion equals
the corresponding row of $V_0$, by its exact A agreement.  Thus the
local decoder recovers all four reference parts and all slots
jointly.  The same good-neighbor count is strictly below $d/2$, so
control majorities are also correct.  This proves the support
containment for every assignment of values on $S$.

For either forward or inverse averaging operator $P$ on $\Gamma$,
the mean-zero bound gives
\[
 \E_k(P1_S(k)-u)^2
 \le\lambda^2u(1-u)\le\lambda^2u.
\]
When $u<\delta_1/4$, Chebyshev's inequality at threshold
$\delta_1/2-u>0$ and a union bound for the two lists prove the first
inequality in \eqref{oc:uncertainty-shrink}.  Its right side is then
at most $32\lambda^2u/\delta_1^2\le u/2$ by
\eqref{oc:lambda-choice}.

Iterate the deterministic operation on supports in
\eqref{oc:uncertainty-set}.  The preceding reasoning holds for every
completion at each step, so the same iterated supports contain all
possible errors.  A depth-$q$ dependency cone is obtained by
composing at most $q$ forward or inverse A moves, with the finite
part and slot bookkeeping.  For fixed $q$ this is a fixed finite
list of permutations.  A consistent assignment to repeated initial
addresses in one cone extends to a global initial word agreeing
with $V_0$ outside $S$.  Applying the global support statement to
that extension proves the claimed result for that cone.  This
argument can use a different global extension for each target.

Assume now $|S|/|\Gamma|<\delta_1/4$.  Every iterated support stays
in this regime, so its integer cardinality decreases by a factor of
at least two while nonempty.  More than $\log_2|\Gamma|$ rounds make
it zero.  Running these rounds synchronously on the complete array
uses $O(|\Gamma|\log|\Gamma|)$ local calls, each of polynomial
bit complexity as established above.  This proves polynomial-time
global correction without using the existential greedy repair in
Lemma~\ref{oc:full-proximity} as an algorithm.
\end{proof}
